% J3.1: el ciclo de optimización responsable.
\documentclass[tikz,border=4pt]{standalone}
\usepackage{resonancia}
\begin{document}
\begin{tikzpicture}
  \node[cajaAzul, minimum width=2.6cm] (m) at (0,0) {\textbf{1. Medir}\\[-1pt]{\footnotesize \texttt{@btime}}};
  \node[cajaVerde, minimum width=2.6cm] (c) at (3.7,0) {\textbf{2. Cuello de botella}\\[-1pt]{\footnotesize ¿dónde se va el tiempo?}};
  \node[cajaNaranja, minimum width=2.6cm] (u) at (7.4,0) {\textbf{3. Cambiar UNA cosa}\\[-1pt]{\footnotesize vistas, tipos, preasignar}};
  \node[caja, draw=rVioleta, fill=rVioleta!8, minimum width=2.6cm] (v) at (11.1,0) {\textbf{4. Volver a medir}\\[-1pt]{\footnotesize ¿mejora? ¿mismo resultado?}};
  \draw[flecha] (m) -- (c); \draw[flecha] (c) -- (u); \draw[flecha] (u) -- (v);
  \draw[flecha, color=rGris] (v.south) -- ++(0,-0.8) -| node[etiqueta, pos=0.25, below] {si merece la complejidad, se repite; si no, se deja como estaba} (m.south);
\end{tikzpicture}
\end{document}
